%! Permutations and Transposes — Unit 2, Lesson 2.4 — Exit Ticket (Answer Key)
\documentclass[10pt]{article}
\usepackage{linalg-article}
\usepackage{linalg-key}   % pulls in -boxes; do NOT also load -boxes
\newcommand{\TallMath}[1]{$\displaystyle #1\rule[-1.4em]{0pt}{3.2em}$}
\newcommand{\xx}{\mathbf{x}}
\newcommand{\bb}{\mathbf{b}}
\newcommand{\tp}{\mathsf{T}}
\newcommand{\PP}{P}

\begin{document}

\pageheader{Unit 2, Lesson 2.4}{Exit Ticket --- Answer Key}

\vspace{0.15in}

\begin{enumerate}[leftmargin=1.8em, itemsep=13pt]
  \item \textbf{Fix the zero pivot and solve.} Let $A=\begin{bmatrix} 0 & 4 \\ 2 & 6 \end{bmatrix}$ and
        $\bb=\begin{bmatrix} 4 \\ 8 \end{bmatrix}$.
        \begin{enumerate}[label=(\alph*), leftmargin=1.8em, itemsep=7pt]
          \item Write the swap matrix $\PP$ and compute $\PP A$ and $\PP\bb$.
\begin{work}
  \PP &= \begin{bsmallmatrix} 0&1\\1&0 \end{bsmallmatrix},\quad
      \PP A = \begin{bsmallmatrix} 2&6\\0&4 \end{bsmallmatrix},\quad
      \PP\bb = \begin{bsmallmatrix} 8\\4 \end{bsmallmatrix}
\end{work}
          \item Solve $A\xx=\bb$ (solve $\PP A\xx=\PP\bb$).
                \hfill $x=\ans{$1$}$\quad $y=\ans{$1$}$
\begin{work}
  4y &= 4,\ y = 1 \\
  2x+6 &= 8,\ x = 1 \\
  0(1)+4(1) &= 4,\quad 2(1)+6(1) = 8 \ \ \text{\checkmark}
\end{work}
        \end{enumerate}
  \item \textbf{Transpose.} Write $A^{\tp}$ for $A=\begin{bmatrix} 2 & 7 \\ 1 & 5 \end{bmatrix}$, and state
        whether $A$ is symmetric.
\begin{work}
  A^{\tp} &= \begin{bsmallmatrix} 2&1\\7&5 \end{bsmallmatrix} \ne A
\end{work}
        \ansline{Not symmetric --- the off-diagonal entries $7$ and $1$ differ.}
  \item \textbf{Why is the swap safe?} A classmate worries that swapping two rows of $A$ might change the
        answer $\xx$. Explain why it does not.
        \ansline{Swapping two rows just reorders the equations (and, with $\bb$ swapped the same way, keeps
        each equation with its right-hand side). The system asks the exact same questions about $\xx$ ---
        only the order changed --- so the solution is unchanged.}
\end{enumerate}

\end{document}
